\chapter{Social Insurance}\label{ch:b3:social-insurance}

Most of what governments spend on people is insurance: against losing a job,
falling ill, becoming disabled, living longer than one's savings. This chapter asks
why these risks are insured by the state rather than by markets, and measures the
trade-off every such programme faces: the more completely it insures, the more it
changes the behaviour whose cost it covers.

\begin{openingchart}{Extra weeks out of work per extra week of potential
unemployment benefit, five estimates from three studies (bars show a range where
the study gives one; Lalive's estimates are lower bounds).}{Katz and Meyer (1990);
Lalive (2008); Johnston and Mas (2018): abstracts of their working papers (1988,
2006, 2016)}
\begin{tikzpicture}
\begin{axis}[omrecord, width=0.5\linewidth, height=0.36\linewidth, xmin=0, xmax=0.6,
  ymin=-0.6, ymax=4.6, ytick={0,...,4}, y dir=reverse,
  yticklabels={Katz--Meyer (United States), Lalive: men (Austria), Lalive: women
    (Austria), Johnston--Mas: out of work (Missouri), Johnston--Mas: on benefits
    (Missouri)}, yticklabel style={font=\scriptsize, align=right, text width=4.4cm},
  xlabel={Weeks per extra week of entitlement}, grid=none, xmajorgrids]
  \addplot[omDef, only marks, mark=*, error bars/.cd, x dir=both, x explicit]
    table[x=mid, y=pos, x error=err] {figdata/bachelor-3/19-social-insurance/out/duration.csv};
\end{axis}
\end{tikzpicture}
\end{openingchart}

% figdata: duration
A week more of entitlement keeps people out of work between 0.09 and 0.32 weeks
longer in these studies, and on benefits about half a week longer. The response is
real and less than one for one: most of the extra entitlement is paid to people who
would have been out of work anyway, or who are not yet near its end.

\begin{recall}
The Year 2 volume showed how adverse selection can unravel an insurance market
and how moral hazard limits insurance, defined the replacement rate, and set
pay-as-you-go pensions against funded ones in the overlapping-generations model.
The regression discontinuity design is chapter 4's.
\end{recall}

\begin{notation}[Symbols of this chapter]\label{not:b3:social-insurance:symbols}
$b$ is a benefit and $r = b/w$ the replacement rate; $D$ the share of time unemployed;
$\varepsilon$ the elasticity of $D$ with respect to $b$; $\gamma$ relative risk aversion; $p$ a survival
probability.
\end{notation}

\section{Why markets for some risks stay thin}\label{sec:b3:social-insurance:why}

\begin{definition}[Social insurance]\label{def:b3:social-insurance:social-insurance}
\index{social insurance}
\emph{Social insurance} is insurance provided or mandated by the state, with
contributions that do not depend on the individual's own risk, against events such
as unemployment, illness, disability, old age and death.
\end{definition}

Three reasons recur. Adverse selection: an insurer who cannot tell high risks from
low ones prices for the average, the low risks leave, and the market shrinks. Risks
that strike many at once, such as a recession, cannot be pooled by a private insurer
whose own capital falls with them. And some risks, such as being born with a
disability, are realised before anyone could buy cover. Mandating cover answers the
first, a state that can borrow answers the second, and only a programme that
covers people from birth answers the third. None answers moral hazard, which
grows with the cover.

\section{Unemployment insurance}\label{sec:b3:social-insurance:ui}

\begin{definition}[Potential benefit duration]\label{def:b3:social-insurance:benefit-duration}
\index{potential benefit duration}
The \emph{potential benefit duration} is the longest period for which an unemployed
person can draw benefits under the rules.
\end{definition}

A worker moves between two states. Employed, she loses her job at the separation
rate $s$ a year; unemployed, she finds one at the job-finding rate $f$. In a steady
state the flows balance, $s(1 - D) = fD$, so the share of time she spends unemployed is
$D = s/(s + f)$. Benefits work on $f$: by lowering the cost of searching less hard or
of turning offers down, they lower the rate at which spells end and raise $D$.

\begin{omfigure}
\begin{tikzpicture}[econ, font=\footnotesize]
  \node[draw, rounded corners, minimum width=2.4cm, minimum height=0.9cm] (E) at (0,0) {employed, $1 - D$};
  \node[draw, rounded corners, minimum width=2.4cm, minimum height=0.9cm] (U) at (5,0) {unemployed, $D$};
  \draw[->, thick, omDef] (E.north east) to[bend left=25] node[above] {separation, $s$} (U.north west);
  \draw[->, thick, omThm] (U.south west) to[bend left=25] node[below] {job finding, $f$; benefits lower it} (E.south east);
\end{tikzpicture}
\par\smallskip{\small The two states of a worker. In a steady state the flows are equal,
$s(1 - D) = fD$, so $D = s/(s + f)$: a lower job-finding rate raises the share of time
unemployed.}
\end{omfigure}

\begin{definition}[Baily--Chetty formula]\label{def:b3:social-insurance:baily-chetty}
\index{Baily--Chetty formula}
The \emph{Baily--Chetty formula} states the condition for the optimal unemployment
benefit: the gain from moving one more dalen to the unemployed, measured by the gap
in marginal utility between unemployment and work, equals the cost to the fund from
the longer spells the dalen causes.
\end{definition}

\begin{proposition}[The Baily--Chetty condition]\label{prop:b3:social-insurance:baily-chetty}
A worker is unemployed a share $D$ of the time, which she chooses by her search, and
receives $b$ when unemployed; a balanced fund levies $t = bD/(1 - D)$ when she works. At
the optimal $b$, $\big(u'(c_u) - u'(c_e)\big)/u'(c_e) = \varepsilon/(1 - D)$. With relative risk aversion $\gamma$ and a consumption
drop $\Delta c/c$, approximately $\gamma\,\Delta c/c = \varepsilon/(1 - D)$.
\end{proposition}

\begin{proof}
Welfare is $W = (1 - D)u(w - t) + Du(b) - \psi(D)$, with $\psi$ the cost of search. Because the worker
chose $D$, the envelope theorem removes $dD/db$ from her own utility, and
$dW/db = Du'(c_u) - (1 - D)u'(c_e)\,dt/db$. From the budget, $(1 - D)\,dt/db = D + (b/(1 - D))\,dD/db = D\big(1 + \varepsilon/(1 - D)\big)$. Setting
$dW/db = 0$ and dividing by $Du'(c_e)$ gives the condition. A first-order expansion of $u'$ gives
$\big(u'(c_u) - u'(c_e)\big)/u'(c_e) \approx \gamma\,\Delta c/c$.
\end{proof}

The formula needs three numbers, all measurable: how much consumption falls on
losing a job, how risk-averse people are, and how much benefits lengthen spells.
It says nothing of the model's details --- what search costs, how wages are set ---
because the measured elasticity summarises them.

\begin{record}[Consumption on losing a job]
Gruber followed the food spending of households in the United States from 1968 to
1987 and matched each spell of unemployment with the benefit the household could
claim under its state's rules. Food spending fell by 7~\% on average when a head of
household became unemployed, and a replacement rate higher by 10 points made the
fall 2.7~\% smaller; without benefits, the fall would have been more than three
times as large. % ledger: b3si:gru.drop
A survey of the labour-supply evidence puts the elasticity of time out of work with
respect to the benefit close to 1.0. % ledger: b3si:km.elast
\end{record}

Put the numbers in the condition. With a drop of 7~\% and relative risk aversion of 2
(an assumption; the problem below uses 3), the left side is $2 \times 0.07 = 0.14$. With $\varepsilon = 1.0$ and $D$ near 5~\%, the
right side is $1.0/0.95 = 1.05$. Taken at face value, the gain from a dalen more of
benefits is a seventh of its cost, and benefits are above the optimum unless risk
aversion is about 15. Three things soften the verdict: food is a necessity, so the
fall in total consumption is larger than the fall in food; part of $\varepsilon$ is not moral
hazard, as the Austrian severance evidence below shows, and that part is a gain;
and the condition holds at the margin, around current benefits. The calculation
shows which number decides the answer: here it is the elasticity.

\begin{omfigure}
\begin{tikzpicture}[econ, x=0.17cm, y=2.4cm]
  \draw[econ axes] (0,0) -- (42,0) node[right, font=\scriptsize] {weeks unemployed};
  \draw[econ axes] (0,0) -- (0,1.25) node[above, font=\scriptsize] {exit rate};
  \draw[curve, omDef] plot[domain=1:24.5, samples=30] (\x, {0.55*exp(-\x/30)+0.12})
    -- (25,0.95) -- (26.5,0.30) -- plot[domain=26.5:40, samples=10] (\x, {0.18+0.1*exp(-(\x-26.5)/6)});
  \draw[curve, omThm, dashed] plot[domain=1:40, samples=40] (\x, {0.50*exp(-\x/30)+0.14});
  \draw[guide] (25,0) node[below, font=\scriptsize] {benefits end} -- (25,1.0);
  \node[font=\scriptsize, omDef, anchor=west] at (27,0.85) {recipients};
  \node[font=\scriptsize, omThm, anchor=west] at (31,0.47) {non-recipients};
\end{tikzpicture}
\par\smallskip{\small A sketch of the prediction, not data: the rate at which recipients
leave unemployment jumps in the weeks benefits run out, while that of people who
never drew benefits does not. The record in the model box below found this pattern
in United States administrative data.}
\end{omfigure}

\begin{model}{Unemployment insurance and search}
\assumes The unemployed choose how hard to search and which offers to accept;
benefits lower the cost of staying unemployed.
\predicts Longer or higher benefits lengthen spells, and the rate at which people
find work jumps when benefits run out.
\tested In United States administrative data, the rate of leaving unemployment
jumps for recipients around the time benefits lapse, but not for nonrecipients;
a week more of entitlement adds 0.16 to 0.20 weeks to a spell. % ledger: b3si:km
In Germany, where entitlement depends on exact age, discontinuities over twenty
years show that the non-employment effect of a month more of benefits is at most
slightly smaller in recessions than in booms, while many more people exhaust their
benefits in recessions. % ledger: b3si:svb
\breaks Part of the response is not moral hazard. In Austria, a severance payment
of two months' earnings, which does not depend on staying unemployed, cut the rate
of finding a job by 8--12~\%, and extending benefits from 20 to 30 weeks cut it by
5--9~\%: people without savings search harder than they would like to, and benefits
let them take longer. % ledger: b3si:ccw
That share of the response is a gain, not a cost, and the formula must count it.
\end{model}

\begin{measure}[The replacement rate]
The replacement rate is the benefit as a share of the previous wage. Rules state
a gross rate, before caps; what matters to behaviour is the
net rate, after the taxes on wages and benefits, often capped at a maximum benefit
and differing with family type and with the length of the spell. Two countries with
the same statutory rate can replace very different shares of what a worker took
home.
\end{measure}

\section{Health insurance: the randomised evidence}\label{sec:b3:social-insurance:health}

\begin{record}[The RAND Health Insurance Experiment, 1974--1981]
More than 5{,}800 people from about 2{,}000 households in six locations in the United
States were randomly assigned to insurance plans that differed in what share of the
bill they paid. Spending fell as the share rose; the experiment's summary elasticity
of medical spending with respect to its out-of-pocket price, $-0.2$, is still the
benchmark, though a single elasticity summarises plans with deductibles and caps
poorly. % ledger: b3si:rand
\end{record}

\begin{record}[The Oregon lottery, 2008]
Oregon drew 30{,}000 names from 90{,}000 low-income adults who had applied, to decide
who could apply for public insurance. In the first year, those drawn were about 25
percentage points more likely to be insured; they used more care, had lower
out-of-pocket spending and medical debt, and reported better physical and mental
health. % ledger: b3si:oregon1
Measured physical health --- blood pressure, cholesterol, blood sugar --- did not
change significantly over about two years, while diabetes diagnoses rose and
depression and financial hardship fell. % ledger: b3si:oregon2
\end{record}

Both experiments find moral hazard in the narrow sense --- people use more care
when it is cheaper --- and both find the main benefit of insurance where insurance
theory puts it: in financial protection, not in measured health over a few years.

\section{Pensions, saving and ageing}\label{sec:b3:social-insurance:pensions}

\begin{definition}[Annuity]\label{def:b3:social-insurance:annuity}
\index{annuity}
An \emph{annuity} is a contract that pays a fixed income for as long as the holder
lives, in exchange for a sum paid in advance.
\end{definition}

\begin{proposition}[Full annuitisation]\label{prop:b3:social-insurance:yaari}
A retiree with wealth $W$ lives one period for sure and a second with probability $p$,
has no bequest motive, and can buy an actuarially fair \omterm{def:b3:social-insurance:annuity}{annuity} or a bond, both at
interest rate zero. If she is risk-averse, she puts all her savings in the \omterm{def:b3:social-insurance:annuity}{annuity}.
\end{proposition}

\begin{proof}
A dalen in the bond pays 1 in period 2 whether she lives or not; a dalen in the fair
\omterm{def:b3:social-insurance:annuity}{annuity} pays $1/p$ if she lives and nothing if she dies. Without a bequest motive,
the payment at death is worth nothing to her. Her budget with the \omterm{def:b3:social-insurance:annuity}{annuity} is
$c_1 + pc_2 = W$; with the bond, $c_1 + c_2 = W$. The \omterm{def:b3:social-insurance:annuity}{annuity}'s budget set contains the bond's,
so she does at least as well with it, and strictly better whenever she saves: the
\omterm{def:b3:social-insurance:annuity}{annuity} raises consumption in the only state in which she consumes.
\end{proof}

Few people buy annuities voluntarily: private annuities are priced for those who
expect to live long (adverse selection again), and public pensions already pay
one. A pay-as-you-go pension is itself an \omterm{def:b3:social-insurance:annuity}{annuity}, financed by current workers,
whose return is the growth of the wage bill: population growth plus productivity
growth.

\begin{omchart}{Old-age dependency ratio (population aged 65 and over per 100 aged
15--64), six economies, 1960--2025.}{World Bank, World Development Indicators
(SP.POP.DPND.OL)}
\begin{tikzpicture}
\begin{axis}[omtime, width=0.66\linewidth, height=0.4\linewidth, xmin=1958, xmax=2027,
  ymin=0, ymax=55, ylabel={Per 100 of working age}, legend pos=north west,
  legend style={font=\footnotesize, legend columns=2}]
  \addplot[omDef, thick, mark=none] table[x=year, y=JPN] {figdata/bachelor-3/19-social-insurance/out/dependency.csv};
  \addplot[omThm, thick, mark=none] table[x=year, y=ITA] {figdata/bachelor-3/19-social-insurance/out/dependency.csv};
  \addplot[black, thick, mark=none] table[x=year, y=DEU] {figdata/bachelor-3/19-social-insurance/out/dependency.csv};
  \addplot[black, thick, dashed, mark=none] table[x=year, y=USA] {figdata/bachelor-3/19-social-insurance/out/dependency.csv};
  \addplot[omProp, thick, mark=none] table[x=year, y=KOR] {figdata/bachelor-3/19-social-insurance/out/dependency.csv};
  \addplot[omDef, thick, dashed, mark=none] table[x=year, y=CHN] {figdata/bachelor-3/19-social-insurance/out/dependency.csv};
  \legend{Japan, Italy, Germany, United States, Korea, China}
\end{axis}
\end{tikzpicture}
\end{omchart}

% figdata: dependency_points
Japan had 9 people of 65 and over per 100 of working age in 1960 and 51 in 2025;
Korea, 6 and 29. A pay-as-you-go pension whose benefits are fixed as a share of
wages needs a contribution rate proportional to this ratio: Japan's has risen
nearly sixfold. Raising the retirement age, cutting replacement rates and building
funded savings are the ways out, and every ageing country uses some mix of them.

\begin{definition}[Crowd-out]\label{def:b3:social-insurance:crowd-out}
\index{crowd-out}
\emph{Crowd-out} is the reduction in private saving or private insurance caused by
a public programme that provides the same protection.
\end{definition}

\begin{record}[Crowd-out measured]
When Medicaid, the United States' public health insurance for low-income families,
was extended between 1987 and 1992, the eligibility of children rose by half and
that of pregnant women doubled. Cutler and Gruber estimated that between 50~\% and
75~\% of the rise in Medicaid coverage was matched by a fall in private coverage:
most of the new public insurance replaced insurance people already had.
% ledger: b3si:cg.crowdout
In Denmark, using 41 million observations on savings, Chetty and co-authors found
that each dollar's worth of government spending on subsidies to retirement accounts
raised total saving by about 1 cent: only some 15~\% of people, the active savers,
respond to subsidies at all, and the other 85~\% do not. Contributions made
automatically raised saving more.
% ledger: b3si:cflno.denmark
\end{record}

\omterm{def:b3:social-insurance:crowd-out}{Crowd-out} decides what a programme adds. A subsidy that buys saving people would
have done anyway, or insurance they already held, costs its full price and adds
little; the same money spent on people who had no protection adds all of it.

\section{Disability insurance}\label{sec:b3:social-insurance:disability}

\begin{definition}[Disability insurance]\label{def:b3:social-insurance:disability}
\index{disability insurance}
\emph{Disability insurance} pays an income to people whose health keeps them from
working, after a test of their capacity to work.
\end{definition}

The test is the programme's hardest part: health is not observed as precisely as
unemployment, and the line between those who cannot work and those for whom
working does not pay moves with the labour market.

\begin{record}[The United States disability rolls after 1984]
Between 1984 and 2000, the share of adults below retirement age receiving federal
disability benefits rose from 3.1~\% to 5.3~\%. Screening had been loosened in 1984, and
benefits replaced a rising share of the falling wages of the less skilled; after
1984, applications became two to three times as responsive to local falls in labour
demand, and less-skilled workers who lost their jobs left the labour force
instead of becoming unemployed. Without the change, the measured unemployment
rate would have been about two-thirds of a point higher. % ledger: b3si:ad
\end{record}

\section{Exercises}\label{sec:b3:social-insurance:exercises}

\begin{exercise}[$\star$]\label{exo:b3:social-insurance:1}
In the opening chart, what is the range of estimates? Why might the effect on time
on benefits exceed the effect on time out of work?
\end{exercise}

\begin{exercise}[$\star$]\label{exo:b3:social-insurance:2}
A worker earned 2{,}500 dalens a month. The benefit replaces 60~\% of the wage up to a
maximum of 1{,}800. Find her benefit and her replacement rate; then those of a worker
who earned 4{,}000.
\end{exercise}

\begin{exercise}[$\star$]\label{exo:b3:social-insurance:3}
Population grows by 0.5~\% a year and productivity by 1.5~\%. What return does a
pay-as-you-go pension pay? What does it pay if population starts to shrink by 0.5~\%?
\end{exercise}

\begin{exercise}[$\star$]\label{exo:b3:social-insurance:4}
What evidence would contradict the claim that unemployment benefits lengthen spells?
\end{exercise}

\begin{exercise}[$\star\star$]\label{exo:b3:social-insurance:5}
Relative risk aversion is 2, the elasticity of unemployment duration with respect to
the benefit is 0.5, and the unemployment rate is 5~\%. Consumption falls by 10~\% on
losing a job. Should benefits rise or fall? What consumption drop would make them
optimal?
\end{exercise}

\begin{exercise}[$\star\star$]\label{exo:b3:social-insurance:6}
Potential duration is extended from 26 to 39 weeks. Using the range of the Katz--Meyer
estimate, how much longer are the spells of recipients?
\end{exercise}

\begin{exercise}[$\star\star$]\label{exo:b3:social-insurance:7}
A retiree survives to the second period with probability 0.8 and the interest rate is
3~\%. What is the fair price of an \omterm{def:b3:social-insurance:annuity}{annuity} paying 1 dalen in period 2? What does a bond
paying 1 dalen cost?
\end{exercise}

\begin{exercise}[$\star\star$]\label{exo:b3:social-insurance:8}
With the WDI data, find the five countries whose old-age dependency ratio rose most
between 2000 and 2025.
\end{exercise}

\begin{exercise}[$\star\star$]\label{exo:b3:social-insurance:9}
What would the Oregon data have shown if insurance improved physical health within
two years?
\end{exercise}

\begin{exercise}[$\star\star\star$]\label{exo:b3:social-insurance:10}
Prove the Baily--Chetty condition, and show where the envelope theorem enters. Why
does the condition hold whatever the cost of search?
\end{exercise}

\begin{exercise}[$\star\star\star$]\label{exo:b3:social-insurance:11}
Half of the retirees survive to period 2 with probability 0.5, half with 0.9; each
knows her type and the insurer does not. Each values an \omterm{def:b3:social-insurance:annuity}{annuity} at her own fair
price and buys only if the price does not exceed it. Show that a pooled fair price
unravels, and find the price at which the market settles.
\end{exercise}

\begin{exercise}[$\star\star\star$]\label{exo:b3:social-insurance:12}
Simulate 20{,}000 job losers aged 40 to 60 whose potential duration is 12 months below 50
and 18 months from 50, with non-employment duration rising by 0.15 months per month of
entitlement, and estimate the jump at 50 with \omterm{def:b3:regression-discontinuity:local-linear}{local linear regressions}. Recover the
per-month effect.
\end{exercise}

\section{Problem: Brenland's Unemployment Fund}\label{sec:b3:social-insurance:problem}

\begin{problem}[{Weekend problem --- Brenland's unemployment fund}]\label{pb:b3:social-insurance:fund}
Brenland's fund pays 60~\% of the previous wage and is financed by a contribution on
wages; take the mean wage as 31{,}000 dalens a year and the 2025 labour-force data of
\code{load("accounts")}. A survey finds that consumption falls on becoming
unemployed by $0.24 - 0.28r$, with $r$ the replacement rate. Relative risk aversion is 3, and
the elasticity of unemployment duration with respect to the benefit is 0.4.
% figdata: fund_stats, fund_curve

\medskip\noindent\textbf{Part I --- the fund.}
\begin{enumerate}
\item How many people were unemployed in 2025, and at what rate?
\item Compute the annual benefit at 60~\%.
\item Compute the fund's spending, treating every unemployed person as unemployed
all year.
\item Find the contribution rate that balances the fund.
\item Compute the consumption drop at 60~\% and with no insurance.
\item Why do the unemployed not smooth consumption on their own savings?
\end{enumerate}

\medskip\noindent\textbf{Part II --- search.}
\begin{enumerate}[resume]
\item By how much does unemployment duration change if $r$ falls from 60~\% to 35~\%?
\item If the unemployment rate moves in proportion, what does it become?
\item What happens to the contribution rate, and why is this a fiscal externality?
\item Give two reasons why higher benefits lengthen spells.
\item What does the Austrian severance-pay result say about these reasons?
\item Why does the distinction matter for welfare?
\end{enumerate}

\medskip\noindent\textbf{Part III --- the optimum.}
\begin{enumerate}[resume]
\item Compute the left side of the approximate Baily--Chetty condition at 60~\%.
\item Compute the right side.
\item Are benefits too generous or too low at 60~\%?
\item Find the optimal replacement rate, holding the unemployment share at 6~\% in
$\varepsilon/(1 - D)$ (at 4.8~\% the right side changes by 1.3~\%).
\item Compute the contribution rate at that rate, with the unemployment rate of
question 8. How much of the fall from question 4 comes from lower unemployment?
\item Compute the consumption drop at that rate.
\item Solve the exact condition with $u'(c) = c^{-3}$. Why is the answer higher?
\item Redo question 16 with an elasticity of 0.2.
\item Redo it with relative risk aversion 2.
\end{enumerate}

\medskip\noindent\textbf{Part IV --- duration.} A regression discontinuity at age 50,
where potential duration rises by 6 months, finds that non-employment lengthens by
0.9 months.
\begin{enumerate}[resume]
\item Compute the effect per month of entitlement.
\item Compare it with the opening chart.
\item What makes the age threshold credible, and what limits the estimate?
\item State the result: Brenland's optimal replacement rate under the measured
elasticity.
\end{enumerate}
\end{problem}
